\documentclass[11pt,a4paper]{report}

\usepackage[utf8]{inputenc}
\usepackage[T1]{fontenc}
\usepackage{amsmath}
\usepackage{amssymb}
\usepackage{array}
\usepackage{booktabs}
\usepackage{geometry}
\usepackage{longtable}
\usepackage{hyperref}
\usepackage{xcolor}

\geometry{margin=1in}
\hypersetup{
  colorlinks=true,
  linkcolor=blue,
  urlcolor=blue,
  citecolor=blue
}
\numberwithin{equation}{chapter}
\setlength{\parskip}{0.5em}
\setlength{\parindent}{0pt}

\newcommand{\code}[1]{\texttt{\detokenize{#1}}}
\newcommand{\rg}{r_g}
\newcommand{\dd}{\mathrm{d}}
\newcommand{\GeV}{\mathrm{GeV}}
\newcommand{\cm}{\mathrm{cm}}
\newcommand{\gcm}{\mathrm{g\,cm^{-3}}}
\newcommand{\sigmanun}{\sigma_{\nu N}}
\newcommand{\SoftwareVersion}{0.9.0}
\newcommand{\PhysicsVersion}{1.3}
\newcommand{\TheoryVersion}{1.3}
\newcommand{\TheoryCompatibleCommit}{8395440}
\newcommand{\TheoryGenerationDate}{2026-07-13}
\newcommand{\TheoryPipelineVersion}{H3-W11}
\newcommand{\TheoryImplementationStatus}{implemented through H3-W10 PYTHIA 8 and HepMC3 Event Generation}

\begin{document}

\begin{titlepage}
\centering
\vspace*{2.0cm}
{\Huge\bfseries HADROS3: Physics and Statistical Theory of the Implemented Pipeline\par}
\vspace{1.0cm}
{\Large Rafael C. R. de Lima\par}
\vfill
{\large Implemented theory through H3-W10 PYTHIA 8 and HepMC3 Event Generation\par}
\vspace{0.5cm}
{\large \TheoryGenerationDate\par}
\vfill
\begin{minipage}{0.82\textwidth}
\small
This document describes only the physics, statistics, numerical methods, and
proxy models implemented in the current HADROS3 pipeline. It is not a user
manual, dashboard guide, installation guide, or proposal for future stages.
Where the implementation uses diagnostic renderers or physics proxies, that
status is stated explicitly.

This is a living theory document for HADROS3. It describes exclusively the
physics currently implemented in the official HADROS3 pipeline. Future or
planned models are not described here until they become part of the official
implementation.
\end{minipage}
\end{titlepage}

\chapter*{Official Physics Reference}
\addcontentsline{toc}{chapter}{Official Physics Reference}

This document is the official reference describing the physics currently
implemented in HADROS3.

If any discrepancy is found between the implementation and this document, either:
\begin{itemize}
\item the implementation must be corrected; or
\item this document must be updated.
\end{itemize}
Such discrepancies must never be ignored.

The scientific traceability chain for implemented physics is
\begin{equation}
\mathrm{code}
\rightarrow
\mathrm{provenance}
\rightarrow
\mathrm{Theory\ PDF}.
\end{equation}
Each simulation provenance record must identify the theory version that
describes its physical and statistical assumptions.

\clearpage
\chapter*{Theory Version Information}
\addcontentsline{toc}{chapter}{Theory Version Information}

\begin{longtable}{>{\raggedright\arraybackslash}p{0.34\textwidth}
                  >{\raggedright\arraybackslash}p{0.52\textwidth}}
\toprule
Field & Value \\
\midrule
\endfirsthead
\toprule
Field & Value \\
\midrule
\endhead
Software Version & \SoftwareVersion \\
Physics Version & \PhysicsVersion \\
Theory Version & \TheoryVersion \\
Compatible HADROS3 commit & \TheoryCompatibleCommit \\
Document generation date & \TheoryGenerationDate \\
Pipeline version & \TheoryPipelineVersion \\
Current implementation status & \TheoryImplementationStatus \\
\bottomrule
\end{longtable}

\begin{longtable}{>{\raggedright\arraybackslash}p{0.35\textwidth}
                  >{\raggedright\arraybackslash}p{0.45\textwidth}}
\toprule
Stage & Theory status \\
\midrule
\endfirsthead
\toprule
Stage & Theory status \\
\midrule
\endhead
H3-W5 UHE Source & implemented and documented \\
H3-W6 Forward Geodesics & implemented and documented \\
H3-W7 DIS Interaction Sampler & implemented and documented \\
H3-W8 Observer Bridge & implemented/proxy and documented \\
H3-W8b Observer Image Branch Analysis & implemented/proxy and documented \\
H3-W9a POWHEG Dry Run & implemented/dry run and documented \\
H3-W9b POWHEG Real Smoke/Free & implemented/smoke+free and documented \\
H3-W10 PYTHIA 8 + HepMC3 Event Generation & implemented and documented \\
H3-W11 GEANT4 & implemented for supported-energy validation; UHE guarded \\
H3-W12 Photon Transport & planned, not documented as implemented \\
H3-W13 Spectra & planned, not documented as implemented \\
\bottomrule
\end{longtable}

\clearpage
\tableofcontents

\chapter*{Scope and Notation}
\addcontentsline{toc}{chapter}{Scope and Notation}

The implemented HADROS3 pipeline is organized as a staged calculation,
\begin{equation}
\mathrm{inputs}\rightarrow \mathrm{run}()\rightarrow
\mathrm{outputs}\rightarrow \mathrm{diagnostics}\rightarrow
\mathrm{provenance}.
\end{equation}
Each implemented stage consumes official products from earlier stages and
writes to its own output directory. The stages covered here are the camera
preview, Kerr geometry utilities, analytic medium, UHE source, forward null
geodesics, DIS optical-depth sampler, Observer Bridge scoring, and POWHEG
dry-run/real hard-process generation, and PYTHIA~8 event generation with
HepMC3 output.

This document is maintained as a required living reference for implemented
physics. When a new HADROS3 stage changes or adds physical equations,
statistical weights, physical approximations, or proxy models, this source file
and the regenerated PDF must be updated with the implementation.

\begin{longtable}{>{\raggedright\arraybackslash}p{0.28\textwidth}
                  >{\raggedright\arraybackslash}p{0.32\textwidth}
                  >{\raggedright\arraybackslash}p{0.25\textwidth}}
\caption{Implemented-stage coverage in this theory document. Planned stages are
listed only to mark that their physics is not yet documented as implemented.}\\
\toprule
HADROS3 stage & Implementation status & Theory documented \\
\midrule
\endfirsthead
\toprule
HADROS3 stage & Implementation status & Theory documented \\
\midrule
\endhead
H3-W5 UHE Source & implemented & yes \\
H3-W6 Forward Geodesics & implemented & yes \\
H3-W7 DIS Interaction Sampler & implemented & yes \\
H3-W8 Observer Bridge & implemented/proxy & yes \\
H3-W8b Observer Image Branch Analysis & implemented/proxy & yes \\
H3-W9a POWHEG Dry Run & implemented/dry run & yes \\
H3-W9b POWHEG Real Smoke/Free & implemented/smoke+free & yes \\
H3-W10 PYTHIA 8 + HepMC3 Event Generation & implemented & yes \\
H3-W11 GEANT4 & implemented/supported-energy only & yes \\
H3-W12 Photon Transport & planned & no \\
H3-W13 Spectra & planned & no \\
\bottomrule
\end{longtable}

All radii in the dynamical Kerr calculation are expressed in gravitational
radius units,
\begin{equation}
\rg = \frac{G M}{c^2}.
\end{equation}
The implementation uses Boyer-Lindquist coordinates
$x^\mu=(t,r,\theta,\phi)$ and the metric signature $(-,+,+,+)$.
The dimensionless Kerr spin is denoted by $a$ in units of $\rg$.
The code clamps spin-dependent operations away from the extremal singular
limit where needed.

\begin{longtable}{>{\raggedright\arraybackslash}p{0.22\textwidth}
                  >{\raggedright\arraybackslash}p{0.67\textwidth}}
\caption{Primary symbols used in the implemented theory.}\\
\toprule
Symbol & Meaning \\
\midrule
\endfirsthead
\toprule
Symbol & Meaning \\
\midrule
\endhead
$r,\theta,\phi$ & Boyer-Lindquist spherical coordinates in $\rg$ units.\\
$\Sigma,\Delta,A$ & Standard Kerr metric functions.\\
$p_\mu$ & Covariant four-momentum components.\\
$E_{\nu,\mathrm{local}}$ & Neutrino energy measured in a local medium or ZAMO frame.\\
$\rho(r,\theta)$ & Analytic torus rest-mass density.\\
$n_b$ & Baryon number density, $n_b=\rho/m_b$.\\
$\sigmanun(E)$ & Tabulated charged-current neutrino-nucleon DIS cross section.\\
$\tau$ & Integrated DIS optical depth along one forward path.\\
$S_{\rm obs}$ & Observer Bridge final observation score.\\
$x$ & Bjorken momentum fraction of the struck parton.\\
$y$ & DIS inelasticity, $y=P\cdot q/(P\cdot k)$.\\
$Q^2$ & Virtuality of the exchanged boson, $Q^2=-q^2>0$.\\
$W$ & Invariant mass of the produced hadronic system.\\
$\sqrt{s}$ & Center-of-mass energy of the neutrino--nucleon system.\\
$G_F$ & Fermi constant.\\
$m_W$ & $W$ boson mass.\\
$V_{q_iq_j}$ & CKM matrix element for the quark flavor transition $q_i\to q_j$.\\
$F_2,\,xF_3,\,F_L$ & CC neutrino DIS structure functions.\\
$Y_\pm$ & DIS angular factors, $Y_\pm = 1\pm(1-y)^2$.\\
$B(\Phi_n)$ & Born differential cross section on $n$-body phase space.\\
$\bar{B}(\Phi_n)$ & NLO-corrected Born (virtual $+$ integrated real counterterms).\\
$\Delta(\Phi_n,p_T)$ & POWHEG Sudakov form factor.\\
\bottomrule
\end{longtable}

\chapter{Camera / Kerr Ray Tracing}

\section{Implemented camera modes}

The camera preview layer is a visual and diagnostic stage. It is separate from
the DIS interaction physics and from the Observer Bridge scoring. Its main
product is a rendered observer-view image, typically
\code{CameraPreview/hadros3_camera_preview.png}. The implemented backend
selection includes:
\begin{itemize}
\item \code{analytic_geometry_only}: always available; a schematic camera
preview, not a physical Kerr ray-traced image.
\item \code{kerr_like_cuda}: a self-contained HADROS3 CUDA preview executable
using a Kerr-like camera model.
\item \code{full_kerr}: a self-contained HADROS3 CUDA preview executable using
the full-Kerr preview mode when available.
\item \code{legacy_cpu_geodesic_preview}: detected for provenance but not used
as the primary self-contained HADROS3 runtime path.
\end{itemize}

The camera preview can render a torus-like visual object and polar funnels.
This visual object is not the DIS optical-depth calculation. The provenance
marks camera-preview products separately from the shared
\code{MediumRenderer}. In current provenance, the camera preview records
\code{medium_renderer_used = false} because its purpose is visual camera
preview, not the shared diagnostic density-map renderer.

\section{Camera frame}

Observer Bridge camera diagnostics use the same geometric convention for the
camera direction that is also relevant when overlaying candidates on the camera
preview. A camera position is built from
\begin{equation}
\mathbf{x}_{\rm obs}
= r_{\rm obs}
\begin{pmatrix}
\sin i \cos\varphi_{\rm obs}\\
\sin i \sin\varphi_{\rm obs}\\
\cos i
\end{pmatrix},
\end{equation}
where $r_{\rm obs}$ is \code{observer_distance_rg}, $i$ is
\code{observer_inclination_deg}, and $\varphi_{\rm obs}$ is
\code{observer_azimuth_deg}. The forward axis points approximately toward the
origin,
\begin{equation}
\hat{\mathbf{f}} =
-\frac{\mathbf{x}_{\rm obs}}{|\mathbf{x}_{\rm obs}|}.
\end{equation}
The right and up axes are built from a world $z$ axis, with a fallback right
axis if the camera is too close to the polar direction:
\begin{align}
\hat{\mathbf{r}}_{\rm cam} &=
\frac{\hat{\mathbf{f}}\times \hat{\mathbf{z}}}
{|\hat{\mathbf{f}}\times \hat{\mathbf{z}}|},\\
\hat{\mathbf{u}}_{\rm cam} &=
\frac{\hat{\mathbf{r}}_{\rm cam}\times \hat{\mathbf{f}}}
{|\hat{\mathbf{r}}_{\rm cam}\times \hat{\mathbf{f}}|}.
\end{align}

\section{FOV diagnostic and Kerr-matched overlay}

The Observer Bridge geometric camera view and FOV score use a pinhole proxy. It
converts a candidate position
$\mathbf{x}$ into a camera vector
\begin{equation}
\mathbf{d}=\mathbf{x}-\mathbf{x}_{\rm obs}
\end{equation}
and projects it as
\begin{align}
z_{\rm cam} &= \mathbf{d}\cdot \hat{\mathbf{f}},\\
x_{\rm ndc} &=
\frac{\mathbf{d}\cdot \hat{\mathbf{r}}_{\rm cam}}
{z_{\rm cam}\tan(\mathrm{FOV}_x/2)},\\
y_{\rm ndc} &=
\frac{\mathbf{d}\cdot \hat{\mathbf{u}}_{\rm cam}}
{z_{\rm cam}\tan(\mathrm{FOV}_y/2)}.
\end{align}
The field-of-view flag is true when $z_{\rm cam}>0$ and
$|x_{\rm ndc}|\leq 1$, $|y_{\rm ndc}|\leq 1$. This proxy is not the default
mapping used by the camera overlay.

The default overlay records
\code{candidate_overlay_projection_model=kerr_geodesic_pixel_match}. It
integrates backward Kerr camera rays on the configured pixel grid and associates
a candidate with the ray of acceptable minimum spatial approach. Thus the
plotted pixel includes Kerr lensing of the camera ray. This remains a
diagnostic position match, not secondary-particle emission or transport. The
explicit \code{candidate_overlay_mapping=geometric_proxy} option retains the
pinhole projection only as a labeled fallback.

Camera Preview stores raw PNG rows with $+e_\theta$ at the top, which is
south-up for the adopted spherical basis. Every spatial Observer Overview
product applies exactly one vertical display transform,
\code{flip_y_for_north_up}; its background and all overlaid pixel coordinates
use the same transform. The resulting top direction is
$-e_\theta/+z$/north and the bottom direction is $+e_\theta/-z$/south.

\section{Inputs and outputs}

\begin{longtable}{p{0.30\textwidth}p{0.60\textwidth}}
\toprule
Role & Implemented item \\
\midrule
Inputs & Black-hole spin, observer distance, inclination, azimuth, FOV, image
resolution, torus and funnel visual parameters.\\
Outputs & \code{CameraPreview/hadros3_camera_preview.png},
\code{CameraPreview/hadros3_camera_preview_summary.json}.\\
Approximations & Analytic fallback is schematic. CUDA preview is visual camera
rendering; candidate overlays remain pinhole proxies.\\
Provenance & Records backend used, fallback status, whether CUDA was used, and
whether full-Kerr preview mode was requested.\\
\bottomrule
\end{longtable}

\chapter{Black Hole and Kerr Geometry}

\section{Metric functions}

HADROS3 uses the Kerr metric in Boyer-Lindquist coordinates and in $\rg$
units. The standard metric functions are
\begin{align}
\Sigma &= r^2 + a^2\cos^2\theta,\\
\Delta &= r^2 - 2r + a^2,\\
A &= (r^2+a^2)^2 - a^2\Delta\sin^2\theta.
\end{align}
The outer horizon radius is
\begin{equation}
r_+ = 1 + \sqrt{1-a^2}.
\end{equation}

The nonzero covariant metric components implemented for the dynamical
calculation are
\begin{align}
g_{tt} &= -\left(1-\frac{2r}{\Sigma}\right),\\
g_{t\phi} &= -\frac{2ar\sin^2\theta}{\Sigma},\\
g_{rr} &= \frac{\Sigma}{\Delta},\\
g_{\theta\theta} &= \Sigma,\\
g_{\phi\phi} &=
\left(r^2+a^2+\frac{2a^2r\sin^2\theta}{\Sigma}\right)\sin^2\theta
= \frac{A\sin^2\theta}{\Sigma}.
\end{align}
The corresponding inverse components used by the Hamiltonian geodesic solver
are
\begin{align}
g^{tt} &= -\frac{A}{\Sigma\Delta},\\
g^{t\phi} &= -\frac{2ar}{\Sigma\Delta},\\
g^{rr} &= \frac{\Delta}{\Sigma},\\
g^{\theta\theta} &= \frac{1}{\Sigma},\\
g^{\phi\phi} &=
\frac{\Delta-a^2\sin^2\theta}{\Sigma\Delta\sin^2\theta}.
\end{align}

\section{ZAMO quantities}

The implemented ZAMO lapse and frame-dragging angular velocity are
\begin{align}
\alpha &= \sqrt{\frac{\Sigma\Delta}{A}},\\
\omega &= \frac{2ar}{A}.
\end{align}
For a covariant momentum with components $p_t$ and $p_\phi$, the ZAMO energy
proxy used in the C++ tetrad utilities is
\begin{equation}
E_{\rm ZAMO}= -\frac{p_t+\omega p_\phi}{\alpha}.
\end{equation}
The local tetrad initialization maps a normalized local spatial direction
$(n_r,n_\theta,n_\phi)$ into a contravariant null vector and then lowers the
index with $g_{\mu\nu}$. The implementation checks that the resulting null norm
is finite and that the ZAMO energy is positive.

\section{Null invariant}

The null condition is
\begin{equation}
g^{\mu\nu}p_\mu p_\nu = 0.
\end{equation}
Both Python reference utilities and the C++ forward-geodesic backend monitor
this invariant. The calculation also monitors the Killing quantities
\begin{align}
E_\infty &= -p_t,\\
L_z &= p_\phi,
\end{align}
which are conserved because the metric is stationary and axisymmetric.

\chapter{Torus / Medium}

\section{Analytic density model}

The implemented medium model is an analytic torus density field. It is not a
hydrodynamical simulation. The density is evaluated as
\begin{equation}
\rho(r,\theta)
=
\rho_0
\exp\left[-\frac{1}{2}
\left(\frac{r-r_{\rm peak}}{\sigma_r}\right)^2\right]
\exp\left[-\frac{1}{2}
\left(\frac{\theta-\pi/2}{\sigma_\theta}\right)^2\right],
\label{eq:torus_density}
\end{equation}
with the hard radial support
\begin{equation}
\rho(r,\theta)=0
\quad \mathrm{if}\quad
r<r_{\rm inner}\quad\mathrm{or}\quad r>r_{\rm outer}.
\end{equation}
The implemented radial width is
\begin{equation}
\sigma_r = \frac{1}{2}(r_{\rm outer}-r_{\rm inner}),
\end{equation}
and the angular width is
\begin{equation}
\sigma_\theta = \max(\theta_{\rm half},10^{-6}).
\end{equation}
The configuration parameter \code{half_opening_angle_deg} is therefore a
Gaussian angular width, not a hard angular boundary. This interpretation is
recorded in diagnostics as
\code{half_opening_angle_interpretation = gaussian_width_not_boundary}.

Within the hard radial support, the code can apply a configured
\code{density_floor_g_cm3} by taking the maximum of the analytic value and the
floor. Outside the radial shell, the density remains exactly zero.

\section{MediumRenderer}

The shared \code{MediumRenderer} renders the same analytic density model used
by the DIS sampler diagnostics. It records
\begin{align}
\code{medium_renderer_used} &= \code{true},\\
\code{density_model_has_hard_radial_cut} &= \code{true},\\
\code{density_model_theta_profile} &= \code{gaussian},\\
\code{density_model_theta_is_hard_cut} &= \code{false}.
\end{align}
The renderer draws the hard radial cuts and Gaussian angular density levels.
Those angular levels are visual density contours, not physical walls.

\section{Inputs and outputs}

\begin{longtable}{p{0.30\textwidth}p{0.60\textwidth}}
\toprule
Role & Implemented item \\
\midrule
Inputs & \code{rho0_g_cm3}, \code{r_inner_rg}, \code{r_peak_rg},
\code{r_outer_rg}, \code{half_opening_angle_deg}, optional density floor.\\
Used by & DIS density evaluation, DIS density diagnostics, forward geometry
diagnostics, Observer Bridge diagnostic projections.\\
Approximation & Static analytic density field; no GRMHD evolution; no
self-consistent thermal or composition model.\\
Output semantics & Positive density only inside the hard radial shell; Gaussian
tail in $\theta$.\\
\bottomrule
\end{longtable}

\chapter{UHE Source}

\section{Coordinate cone sampling}

The implemented UHE source is a coordinate-volume sampler in a polar cone.
For a cone spanning $r_{\min}\leq r\leq r_{\max}$ and
$\theta_{\min}\leq \theta\leq\theta_{\max}$, the coordinate volume used for
sampling is
\begin{equation}
V_{\rm cone}
= \frac{2\pi}{3}
\left(r_{\max}^3-r_{\min}^3\right)
\left(\cos\theta_{\min}-\cos\theta_{\max}\right).
\end{equation}
Uniform coordinate-volume sampling is implemented by drawing
$u_r,u_\theta,u_\phi\sim U(0,1)$ and setting
\begin{align}
r &= \left[
r_{\min}^3 + u_r(r_{\max}^3-r_{\min}^3)
\right]^{1/3},\\
\cos\theta &=
\cos\theta_{\min}
-u_\theta(\cos\theta_{\min}-\cos\theta_{\max}),\\
\phi &= 2\pi u_\phi.
\end{align}
The source sampling and physical PDFs are both recorded as
\begin{equation}
p_{\rm source} = \frac{1}{V_{\rm cone}},
\end{equation}
so the implemented source weight is unity for the default model:
\begin{equation}
w_{\rm source}=
\frac{p_{\rm physical}}{p_{\rm sampling}}=1.
\end{equation}

\section{Energy and directions}

The implemented source energy model is monoenergetic:
\begin{equation}
E_{\nu,\rm emit}=E_0.
\end{equation}
For the isotropic local direction model, the code samples a unit vector in a
ZAMO orthonormal frame:
\begin{align}
\cos\alpha &= 2u_\alpha-1,\\
\beta &= 2\pi u_\beta,\\
n_r &= \cos\alpha,\\
n_\theta &= \sin\alpha\cos\beta,\\
n_\phi &= \sin\alpha\sin\beta.
\end{align}
The sampling PDF is
\begin{equation}
p_\Omega = \frac{1}{4\pi},
\end{equation}
and the default direction weight is unity.

The source stage does not itself produce the physical Kerr four-momentum for
forward propagation. It records the source position, energy, and local direction.
The forward-geodesic stage constructs the physical Kerr null momentum from
those quantities.

\section{Inputs and outputs}

\begin{longtable}{p{0.30\textwidth}p{0.60\textwidth}}
\toprule
Role & Implemented item \\
\midrule
Inputs & Cone geometry, sample count, random seed, monoenergetic
neutrino energy, direction model.\\
Outputs & \code{UHEsource/uhe_neutrino_source_samples.jsonl} and source summary
products.\\
Weights & \code{source_weight = 1} and \code{direction_weight = 1} for the
implemented default uniform/properly matched sampling.\\
Approximation & Coordinate cone volume, not a GR invariant emissivity integral.\\
\bottomrule
\end{longtable}

\chapter{Forward Geodesics}

\section{Momentum construction}

The forward stage consumes UHE source samples and constructs a Kerr null
covector. For a local isotropic direction, the ZAMO tetrad maps the local
direction into a null four-vector and then into covariant momentum components.
For legacy coordinate radial directions, the implementation constructs a
covariant momentum with $p_t=-E$ and solves the null constraint for $p_r$.

The implementation records that the physical Kerr four-momentum is constructed
in the forward stage, not in the source sampler.

\section{Hamiltonian system}

The Hamiltonian for null geodesics is
\begin{equation}
H(x,p)=\frac{1}{2}g^{\mu\nu}(x)p_\mu p_\nu.
\end{equation}
The implemented equations of motion are
\begin{align}
\frac{\dd x^\mu}{\dd\lambda} &=
g^{\mu\nu}p_\nu,\\
\frac{\dd p_i}{\dd\lambda} &=
-\frac{1}{2}
\partial_i g^{\mu\nu}p_\mu p_\nu,
\qquad i\in\{r,\theta\}.
\end{align}
The cyclic momenta are conserved by construction in the differential system:
\begin{equation}
\frac{\dd p_t}{\dd\lambda}=0,\qquad
\frac{\dd p_\phi}{\dd\lambda}=0.
\end{equation}

The production C++ implementation uses an adaptive Runge--Kutta--Fehlberg
stepper with invariant checks; the Python reference uses RK4 substeps and may
renormalize $p_r$ to keep the null constraint within tolerance. Derivatives of
inverse metric components are evaluated according to the configured backend.
The solver stops when it reaches the horizon neighborhood, the configured outer
radius, maximum steps, or an invalid invariant condition.

\section{Segment quantities}

Each accepted forward path is written as a set of path segments. The adaptive
integrator records the actual accepted affine interval $\Delta\lambda_i$ and
the momentum normalization $E_{\rm norm}$ used to evolve the dimensionless
covector. For medium four-velocity $u^\mu$, the matter-frame path length is
\begin{equation}
\Delta \ell_{{\rm com},i}
=
\frac{-u^\mu p_\mu}{E_{\rm norm}}\,\Delta\lambda_i\,r_g.
\end{equation}
The contraction and cross-section energy are evaluated at the same segment
midpoint. The ratio by $E_{\rm norm}$ accounts for the serialized GeV momentum
versus the dimensionless momentum used by the affine integrator. Consequently,
$p_\mu\rightarrow c p_\mu$ together with
$\Delta\lambda\rightarrow\Delta\lambda/c$ leaves
$\Delta\ell_{\rm com}$ unchanged. The older Boyer--Lindquist spatial-slice
distance remains in the segment record only as a geometric diagnostic and is
not used for optical depth.

\section{Inputs and outputs}

\begin{longtable}{p{0.30\textwidth}p{0.60\textwidth}}
\toprule
Role & Implemented item \\
\midrule
Inputs & \code{UHEsource/uhe_neutrino_source_samples.jsonl},
black-hole spin, step settings, tolerances, outer radius, requested number of
samples.\\
Outputs & \code{ForwardGeodesics/uhe_neutrino_forward_paths.jsonl},
\code{ForwardGeodesics/uhe_neutrino_forward_path_segments.jsonl}, and summary
diagnostics.\\
Invariants & Null norm, Killing energy, and $L_z$ are monitored.\\
Approximation & Midpoint evaluation of the covariant matter-frame path length
and local energy over each finite affine segment.\\
\bottomrule
\end{longtable}

\chapter{DIS Interaction Sampler}

\section{Density and local energy}

The DIS sampler consumes forward path segments and evaluates the analytic torus
density at segment midpoints. The baryon density is
\begin{equation}
n_b(r,\theta)=\frac{\rho(r,\theta)}{m_b},
\end{equation}
where $m_b$ is the baryon mass in grams.

The local neutrino energy is computed from the covariant momentum. If the
configured static medium model is valid at that position, the code uses the
static observer expression
\begin{equation}
E_{\rm local} = -p_t u^t,\qquad
u^t=\frac{1}{\sqrt{-g_{tt}}}.
\end{equation}
If a static observer is not valid, or if the implementation needs the fallback,
the ZAMO expression is used:
\begin{equation}
E_{\rm local}
=-\left(\frac{p_t}{\alpha}+\frac{\omega p_\phi}{\alpha}\right)
=-\frac{p_t+\omega p_\phi}{\alpha}.
\end{equation}
The static-to-ZAMO fallback is recorded as a physics-risk approximation.

\section{Cross-section interpolation}

The DIS charged-current cross sections are read from compact tabulated files
for the implemented GBW and IIM models. If an energy lies inside the tabulated
range, the interpolation is log-log:
\begin{equation}
\ln \sigmanun(E)
=
\ln \sigma_0
+
\frac{\ln E-\ln E_0}{\ln E_1-\ln E_0}
\left(\ln\sigma_1-\ln\sigma_0\right).
\end{equation}
Requests outside the tabulated range are rejected for that segment and the
segment contribution is set to zero in the diagnostic recomputation.

\section{Optical depth and probability}

For each segment,
\begin{equation}
\Delta\tau_i =
n_{b,i}\,\sigmanun(E_{{\rm local},i})\,
\Delta \ell_{{\rm com},i},
\end{equation}
where $E_{{\rm local},i}=-u^\mu p_\mu$ and
$\Delta\ell_{{\rm com},i}$ is converted from $\rg$ to cm. This is the discrete
midpoint approximation to the invariant integral
\begin{equation}
\tau=\int n_b\,\sigmanun(-u\!\cdot\!p)\,
\frac{-u\!\cdot\!p}{E_{\rm norm}}\,\dd\lambda\,r_g.
\end{equation}
The total optical depth along a path is
\begin{equation}
\tau = \sum_i \Delta\tau_i.
\end{equation}
The interaction probability is
\begin{equation}
P_{\rm int}=1-\exp(-\tau),
\end{equation}
implemented numerically as
\begin{equation}
P_{\rm int}=-\operatorname{expm1}(-\tau)
\end{equation}
with clipping to the interval $[0,1]$.

\section{Interaction sampling}

If a path is accepted by the Bernoulli trial with probability $P_{\rm int}$,
the segment index is sampled by the inverse CDF
\begin{equation}
P(i)=\frac{\Delta\tau_i}{\sum_j \Delta\tau_j}.
\end{equation}
The Bernoulli trial is evaluated for every path, independent of
\code{max_interactions}. If the number of preliminary successes exceeds that
cap, HADROS3 assigns each success an independent deterministic hash priority
from the configured seed and event identifier, and retains the smallest
priorities. Conditional on the preliminary successes this is a uniform
subsample without replacement. It is reproducible and invariant under a
permutation of the input path order.

After the segment is selected, let $u\in[0,1]$ linearly interpolate its
coordinates. The midpoint local energy, cross section, and affine conversion
are held fixed within that segment, while the medium density is evaluated at
129 nodes. A 128-cell trapezoidal quadrature constructs the intrasegment CDF
proportional to
\begin{equation}
\frac{\dd\tau_i}{\dd u}\propto
n_b\!\left(r(u),\theta(u)\right)
\sigmanun(E_{{\rm local},i})
\frac{\dd\ell_{{\rm com},i}}{\dd u}.
\end{equation}
The density is linearly interpolated inside each quadrature cell and its local
integral is inverted numerically. Thus the location follows the implemented
optical-depth density rather than a uniform coordinate draw. The sampler also
requires the final recorded interaction point to have positive medium density
and records
\begin{align}
\code{interaction_point_density_checked} &= \code{true},\\
\code{interaction_point_inside_medium} &= \code{true}\quad \mathrm{when}\quad
\rho(r_{\rm int},\theta_{\rm int})>0,\\
\code{interaction_point_rho_g_cm3} &= \rho(r_{\rm int},\theta_{\rm int}).
\end{align}
Rare boundary/fallback behavior is recorded in
\code{interaction_point_sampling_method}. The present criterion for a valid
accepted point is $\rho(r_{\rm int},\theta_{\rm int})>0$.

\section{Weights}

The DIS products preserve source and direction weights from the UHE source.
The pre-event weight used downstream is recorded as
\code{final_pre_event_weight}. The DIS sampler also records optical-depth and
expected-interaction quantities so that later stages can distinguish the
sampling decision from pre-event observation scoring.

\section{Inputs and outputs}

\begin{longtable}{p{0.30\textwidth}p{0.60\textwidth}}
\toprule
Role & Implemented item \\
\midrule
Inputs & Forward geodesic paths and segments, source samples, DIS cross-section
tables, analytic medium parameters, random seed.\\
Outputs & \code{DIS/dis_path_optical_depths.jsonl},
\code{DIS/dis_accepted_interactions.jsonl}, optical-depth report, diagnostics,
GBW/IIM comparison products.\\
Diagnostics & Tau distribution, interaction probability distribution, optical
depth maps, medium density map, accepted interaction distributions, sigma and
correlation plots.\\
Approximation & Static or ZAMO medium four-velocity; midpoint quadrature of
density, local energy, cross section, and covariant comoving path length.\\
\bottomrule
\end{longtable}

\chapter{Observer Bridge}

\section{Purpose}

The Observer Bridge consumes accepted DIS interactions and assigns a
non-destructive observation score. It does not generate secondary particles,
does not call POWHEG, does not call PYTHIA, does not call GEANT4, and does not
perform photon transport. All accepted DIS interactions become Observer Bridge
candidates, even if they are outside the camera FOV or have small score.

\section{Geometry proxies}

The interaction position is converted to Cartesian coordinates through
\begin{equation}
(x,y,z)=
\left(
r\sin\theta\cos\phi,\,
r\sin\theta\sin\phi,\,
r\cos\theta
\right).
\end{equation}
The observer is placed at the configured camera position. The geometric
direction from the interaction to the observer is
\begin{equation}
\hat{\mathbf{n}}_{\rm obs}
=
\frac{\mathbf{x}_{\rm obs}-\mathbf{x}_{\rm int}}
{|\mathbf{x}_{\rm obs}-\mathbf{x}_{\rm int}|}.
\end{equation}
The implemented escape proxy compares this direction with the local radial
direction,
\begin{equation}
w_{\rm esc}
=
\max\left(0,\,
\hat{\mathbf{x}}_{\rm int}\cdot \hat{\mathbf{n}}_{\rm obs}
\right).
\end{equation}
This is a proxy, not a particle-transport escape calculation.

The FOV weight is either a hard FOV flag or a soft Gaussian-like angular proxy,
depending on configuration. Visibility, line of sight, and redshift are not
physical models in H3-W8: all three are immutable unity diagnostic factors,
recorded as \code{unity_diagnostic_not_physics} with their corresponding
\code{*_physics_model=false}. The legacy enable/check flags are retained for
configuration compatibility but do not change these unit weights. In
particular, no metric redshift integral or occultation test is implied.

\section{Score definition}

The implemented physics weight in the C++ Observer Bridge is
\begin{equation}
w_{\rm physics}=\code{final_pre_event_weight},
\end{equation}
with a fallback to source times direction weight when that field is absent.
The observer weight is the product of the implemented proxy factors:
\begin{equation}
w_{\rm observer}
=
w_{\rm esc}\,
w_{\rm visibility}\,
w_{\rm FOV}\,
w_{\rm distance}\,
w_{\rm redshift}\,
w_{\rm LOS}.
\end{equation}
The final observation score is
\begin{equation}
S_{\rm obs}
=
w_{\rm physics}\,w_{\rm observer}.
\end{equation}
All three quantities are recorded for each candidate:
\code{physics_weight}, \code{observer_weight}, and
\code{final_observation_score}. The implementation enforces non-negative proxy
weights by clipping or construction.

\section{Proxy risk flags}

The reports record the proxy nature of the bridge:
\begin{align}
\code{secondary_particle_proxy_model} &\neq \code{full_transport},\\
\code{escape_proxy_physics_risk} &= \code{true},\\
\code{visibility_proxy_physics_risk} &= \code{true},\\
\code{redshift_proxy_physics_risk} &= \code{true}.
\end{align}
The stage status is scoring-only. It is intentionally non-destructive.

\section{Inputs and outputs}

\begin{longtable}{p{0.30\textwidth}p{0.60\textwidth}}
\toprule
Role & Implemented item \\
\midrule
Inputs & \code{DIS/dis_accepted_interactions.jsonl} and runtime configuration.
It may consult previous products for diagnostics.\\
Outputs & \code{ObserverBridge/observer_bridge_candidates.jsonl},
\code{ObserverBridge/observer_bridge_ranked_events.jsonl}, JSON/CSV summaries,
camera-view and overlay diagnostics.\\
Invariant policy & Number of candidates scored equals number of accepted DIS
interactions consumed. FOV is not used as a destructive filter.\\
Approximation & Geometric escape, visibility, FOV, redshift, and line-of-sight
proxies. No secondary-particle generation.\\
\bottomrule
\end{longtable}

\chapter{Observer Image Branch Analysis}

H3-W8b introduces an explicit observational image-branch analysis between the
Observer Bridge and POWHEG. The Observer Bridge remains a non-destructive
candidate scorer: it maps every accepted DIS interaction into an observational
candidate and records a closest Kerr-ray diagnostic. H3-W8b asks a different
question. Instead of assigning
\begin{equation}
\mathrm{candidate}\rightarrow\arg\min_{\rm ray} d_{\rm closest},
\end{equation}
it groups all compatible camera rays into observed image branches and selects a
primary branch by an auditable proxy score.
The dashboard presents this stage visually as \code{Gravitational Image
Analysis}, while the legacy \code{ObserverImageBranches/} directory and base
file names remain the compatibility contract.

\section{Inputs}

The implemented stage consumes official Observer Bridge products:
\begin{align}
&\code{ObserverBridge/observer_candidate_kerr_pixel_map.jsonl},\\
&\code{ObserverBridge/observer_bridge_selected_candidates.jsonl},\\
&\code{ObserverBridge/observer_bridge_ranked_events.jsonl}.
\end{align}
When available, it also consumes the Observer Bridge multi-image audit product
\code{ObserverBridge/candidate_multi_image_audit.jsonl}, which records sampled
camera rays passing within the configured matching radius of each candidate.
No DIS interaction, forward geodesic, camera preview, or Kerr integrator result
is modified by H3-W8b.

\section{Image branch definition}

For a candidate $c$, let $\mathcal{R}_c$ be the set of sampled camera rays whose
minimum spatial separation from the candidate is below the matching tolerance.
Rays are clustered in image-pixel space. Each connected cluster is an image
branch: two rays are adjacent when their Euclidean pixel separation is no
greater than the configured link radius, and a branch is a connected component
of this undirected graph. Consequently the relation is transitive and the
result is invariant under input-ray ordering. Each branch $B_{c,j}$ records:
\begin{equation}
B_{c,j} =
\left\{
N_{\rm rays},
\bar{x}_{\rm pix},
\bar{y}_{\rm pix},
\mathrm{Cov}(x_{\rm pix},y_{\rm pix}),
\bar{d}_{\rm closest},
d_{\rm min},
d_{\rm max},
s_{\rm pix}
\right\}.
\end{equation}
Here $N_{\rm rays}$ is the number of compatible rays, $(\bar{x}_{\rm pix},
\bar{y}_{\rm pix})$ is the branch centroid on the camera preview pixel plane,
and $s_{\rm pix}$ is the RMS pixel spread around the centroid.

\section{Branch score proxy}

The primary observed branch is not selected by the single closest ray. Instead,
H3-W8b assigns each branch a proxy score
\begin{equation}
S_{\rm branch}
=
w_{\rm rays}\,w_{\rm closest}\,w_{\rm compactness},
\end{equation}
with the implemented proxy weights
\begin{align}
w_{\rm rays} &= N_{\rm rays},\\
w_{\rm closest} &= \frac{1}{\max(\bar{d}_{\rm closest},\epsilon)},\\
w_{\rm compactness} &= \frac{1}{\max(s_{\rm pix},1\,{\rm pixel})}.
\end{align}
This score favors image branches supported by many rays, small mean closest
approach, and compact clusters on the image plane. It is a ranking proxy, not a
full magnification, radiative transfer, or detector-response calculation.
Reports and provenance explicitly mark this with
\code{branch_score_is_proxy=true} and
\code{branch_score_not_true_magnification=true}.

\section{Primary branch selection}

For each candidate, H3-W8b selects
\begin{equation}
B_{c,\rm primary}
=
\arg\max_j S_{\rm branch}(B_{c,j}).
\end{equation}
The selected primary branches are written to
\code{ObserverImageBranches/observer_image_primary_branches.jsonl}. Downstream
POWHEG preparation consumes that file, so POWHEG receives one primary observed
image branch per selected Observer Bridge candidate and does not perform its
own image-branch selection.

\section{Outputs and provenance}

The stage writes exclusively under \code{ObserverImageBranches/}:
\code{observer_image_branches.jsonl},
\code{observer_image_primary_branches.jsonl},
\code{observer_image_branch_summary.json},
\code{observer_image_branch_report.json},
\code{observer_image_statistics.json},
diagnostic PNGs, a CSV table, and \code{observer_branch_view.html}. Provenance
records \code{observer_image_branch_analysis_invoked=true},
\code{gravitational_image_analysis_invoked=true},
\code{visual_stage_name=Gravitational Image Analysis},
\code{branch_scoring_model=ray_count_closeness_compactness_proxy},
\code{primary_branch_selection_model=argmax_branch_score},
\code{mean_branches_per_candidate}, \code{fraction_multiple_images}, and
\code{primary_branch_selection_proxy=true}. The stage also writes
\code{gravitational_image_*} explanatory products for multiplicity, score
decomposition, primary selection, and POWHEG forwarding without changing the
selection algorithm.

When no multi-image audit row exists but a closest-pixel map is available, the
compatibility layer may emit one synthetic single-ray branch. Such a record is
always tagged
\code{multiplicity_audited=false} and
\code{branch_is_synthetic_fallback=true}; it is not counted as evidence for a
single-image multiplicity claim. Conversely, an audit row that explicitly
contains zero clusters remains a zero-branch result and is never replaced by
the synthetic fallback. These flags are preserved in the primary downstream
record and in provenance.

\section{Physical limitations}

The current branch score is observationally motivated but approximate. It does
not compute rigorous Kerr caustic magnification, phase-space beam Jacobians,
secondary-particle emission, absorption, or detector response. It should be
treated as a transparent primary-image proxy until a later stage implements a
more complete lensing and transport model.

\chapter{POWHEG Hard Process Preparation and Smoke Execution}

\section{Hard subprocess: charged-current neutrino DIS}

The implemented POWHEG process is \texttt{nudis} from POWHEG-BOX-RES. HADROS3
exposes the perturbative order explicitly. The default
\code{perturbative_order=LO} is a Born-only validation/smoke configuration and
writes \code{LOevents 1}. With \code{perturbative_order=NLO}, that flag is
absent and POWHEG evaluates the real and virtual contributions and generates
the hardest QCD emission. The Born-level hard subprocess is
\begin{equation}
\nu_\ell(k) + q_i(p) \;\longrightarrow\; \ell^-(k') + q_j(p'),
\label{eq:dis_born}
\end{equation}
mediated by a virtual $W^+$ boson with four-momentum transfer
$q^\mu = k^\mu - k'^\mu$. The incoming neutrino carries the full UHE energy
$E_{\nu,\rm local}$ from the accepted DIS interaction point; the target is a
parton $q_i$ drawn from the nucleon PDF with momentum fraction $x$. HADROS3
specifies the two beam particles as
\begin{align}
\texttt{ih1} &= 12 \quad (\nu_e,\; \text{PDG electron neutrino}),\\
\texttt{ih2} &= 1  \quad (\text{proton/nucleon, PDF target}).
\end{align}
The CC process is selected in the POWHEG input card by
\texttt{channel\_type}$=3$. The \texttt{vtype} input controls photon/Z content
for neutral-current channels in this backend and does not select the CC lepton
flavor; it is therefore not used by HADROS3 to infer CC flavor.

\section{DIS kinematic variables}

The standard DIS kinematic invariants are built from the neutrino
four-momentum $k^\mu$, the nucleon four-momentum $P^\mu$, and the
momentum-transfer four-momentum $q^\mu = k^\mu - k'^\mu$:
\begin{align}
s &= (k+P)^2 \approx 2\,m_N\,E_{\nu,\rm local},\label{eq:dis_s}\\
Q^2 &= -q^2 > 0,\label{eq:dis_Q2}\\
x &= \frac{Q^2}{2\,P\cdot q},\label{eq:dis_x}\\
y &= \frac{P\cdot q}{P\cdot k} = \frac{Q^2}{xs},\label{eq:dis_y}\\
W^2 &= (P+q)^2 = m_N^2 + Q^2\frac{1-x}{x}.\label{eq:dis_W2}
\end{align}
The approximation for $s$ in equation~\eqref{eq:dis_s} holds at UHE energies
where $E_{\nu,\rm local}\gg m_N$. The Bjorken variable $x\in(0,1]$ is the
momentum fraction of the struck parton, $y\in(0,1]$ is the inelasticity
(fractional energy transfer to the hadronic system in the nucleon rest frame),
and $W$ is the invariant mass of the produced hadronic system.

\section{CC neutrino--nucleon DIS cross section}

The double-differential cross section for inclusive $\nu_\ell N$ CC DIS is
expressed through three structure functions $F_2$, $F_L$, and $xF_3$:
\begin{equation}
\frac{\dd^2\sigma^{\rm CC}}{\dd x\,\dd y}
=
\frac{G_F^2\,s}{4\pi}
\left(\frac{m_W^2}{m_W^2+Q^2}\right)^{\!2}
\left[
  Y_+\,F_2(x,Q^2)
  -y^2 F_L(x,Q^2)
  +Y_-\,xF_3(x,Q^2)
\right],
\label{eq:dis_xsec}
\end{equation}
where $G_F$ is the Fermi constant, $m_W\simeq 80.4\,\mathrm{GeV}$ is the $W$
boson mass, and the inelasticity factors are
\begin{equation}
Y_\pm = 1 \pm (1-y)^2.
\end{equation}
The sign in front of $xF_3$ is $+Y_-$ for neutrinos (as written) and
$-Y_-$ for antineutrinos. The propagator factor
$(m_W^2/(m_W^2+Q^2))^2$ becomes significant when $Q^2 \gtrsim m_W^2$.

At leading order in QCD the structure functions are expressed through the
parton distributions weighted by CKM elements. For CC $\nu N$ DIS the struck
parton must be a down-type quark $q_i\in\{d,s,b\}$ (or an up-type antiquark
from the sea), and the relevant combinations are
\begin{align}
F_2^{\rm CC}(x,Q^2)
&= 2x\sum_i \left|V_{q_iq_j}\right|^2 f_{q_i}(x,Q^2),
\label{eq:F2}\\
xF_3^{\rm CC}(x,Q^2)
&= 2x\sum_i \left|V_{q_iq_j}\right|^2
   \operatorname{sign}(q_i)\,f_{q_i}(x,Q^2),
\label{eq:xF3}
\end{align}
where $V_{q_iq_j}$ is the CKM element for the transition $q_i\to q_j$ and
$\operatorname{sign}(q_i)=+1$ ($-1$) for quarks (antiquarks).
The longitudinal structure function $F_L = F_2 - 2xF_1$ vanishes at LO
(Callan--Gross relation) but receives $\mathcal{O}(\alpha_s)$ NLO corrections.

\section{The POWHEG method}

In the configured NLO mode, POWHEG generates the first (hardest) QCD emission
with NLO accuracy and provides it as input to a subsequent parton shower. The
starting point is the
NLO-corrected Born weight
\begin{equation}
\bar{B}(\Phi_n) = B(\Phi_n) + V(\Phi_n)
+ \int\!\dd\Phi_{\rm rad}\,C(\Phi_{n+1}),
\label{eq:powheg_bbar}
\end{equation}
where $\Phi_n$ is the $n$-body Born phase space, $B(\Phi_n)$ is the Born
squared amplitude, $V(\Phi_n)$ is the renormalized one-loop virtual
correction, and $C(\Phi_{n+1})$ is the integrated infrared counterterm
that removes double-counting with the real emission. The POWHEG master formula
for the fully-differential cross section is
\begin{equation}
\dd\sigma_{\rm POWHEG}
=
\bar{B}(\Phi_n)\,\dd\Phi_n
\Biggl[
\Delta(\Phi_n,\,p_{T,\min})
+\sum_\alpha
\frac{R_\alpha(\Phi_{n+1})}{B(\Phi_n)}\,
\Delta(\Phi_n,\,k_{T,\alpha})\,
\dd\Phi_{\rm rad,\alpha}
\Biggr],
\label{eq:powheg_master}
\end{equation}
where the sum is over all QCD singular regions $\alpha$, $R_\alpha$ is the
real squared amplitude in region $\alpha$, and $k_{T,\alpha}$ is the emission
transverse momentum in that region. The POWHEG Sudakov form factor is
\begin{equation}
\Delta(\Phi_n,\,p_T)
=
\exp\!\left[
-\sum_\alpha
\int_{p_T}^{\infty}
\frac{R_\alpha(\Phi_{n+1})}{B(\Phi_n)}
\,\dd\Phi_{\rm rad,\alpha}
\right].
\label{eq:powheg_sudakov}
\end{equation}
The Sudakov factor ensures that the hardest emission is generated with the
correct NLO-improved distribution. Events generated by
equation~\eqref{eq:powheg_master} carry predominantly positive weights---the
property that gives POWHEG its name---and reproduce the NLO-accurate inclusive
rate after integration.

For the \texttt{nudis} process the Born phase space $\Phi_n$ covers the
two-body $\nu_\ell q_i \to \ell^- q_j$ kinematics parameterized by
$(x,y,Q^2)$. The NLO real corrections include quark-gluon emission
\begin{equation}
\nu_\ell + q_i \;\longrightarrow\; \ell^- + q_j + g
\end{equation}
and the gluon-splitting channel
\begin{equation}
\nu_\ell + g \;\longrightarrow\; \ell^- + q_j + \bar{q}_i,
\end{equation}
both of which can produce a gluon or an extra $q\bar{q}$ pair among the
outgoing hard-process particles in the LHE record. These terms and channels are
disabled in the explicit LO mode.

After the hardest emission is fixed by the POWHEG Sudakov, the parton shower
(Pythia~8, not yet invoked in the current HADROS3 pipeline) adds further
softer radiation using the POWHEG emission scale as its starting scale.

\section{Center-of-mass energy and scale}

Each accepted DIS interaction vertex in HADROS3 provides the local-frame
neutrino energy $E_{\nu,\rm local}$ to POWHEG as \texttt{ebeam1}. The
card sets \code{fixed_target=1}, so \texttt{ebeam2} is interpreted as the
nucleon mass rather than as the energy of an opposing massless beam. The
invariant mass of the neutrino--nucleon hard system is
\begin{equation}
\sqrt{s} = \sqrt{2\,m_N\,E_{\nu,\rm local}}.
\end{equation}
For the UHE range $E_{\nu,\rm local}\sim 10^9$--$10^{12}\,\mathrm{GeV}$,
this gives $\sqrt{s}\sim 43$--$1400\,\mathrm{TeV}$. The upper kinematic
bound on the virtuality integration is
\begin{equation}
Q_{\max}
=
\min\!\left(\sqrt{2m_N\,E_{\nu,\rm local}},\;10^5\,\mathrm{GeV}\right),
\end{equation}
and the lower bound is fixed at $Q_{\min}=10\,\mathrm{GeV}$. The cap at
$10^5\,\mathrm{GeV}$ prevents integration beyond kinematic regions where the
fixed-order NLO approximation is reliable. Renormalization and factorization
scales are set equal to the virtuality $Q$ at each phase-space point
(\texttt{runningscales}$=1$):
\begin{equation}
\mu_R = \mu_F = Q.
\end{equation}

\section{Parton distributions and CKM flavor structure}

HADROS3 uses NNPDF31\_nlo\_as\_0118 (LHAPDF set 303400) for both incoming
beams. This is perturbatively matched to the NLO mode; in the LO smoke mode it
is retained for controlled pipeline continuity and is explicitly not described
as an order-matched LO prediction. The strong coupling $\alpha_s$ is taken from
the same PDF set at each renormalization scale
(\texttt{alphas\_from\_pdf}$=1$).

For CC $\nu_\ell N$ DIS the incoming parton must be a down-type quark. The
weak transition to an up-type quark is weighted by the squared CKM elements.
The dominant channels and their approximate weights are
\begin{align}
d &\to u: \quad |V_{ud}|^2 \approx 0.947,\\
s &\to c: \quad |V_{cs}|^2 \approx 0.947,\\
b &\to t: \quad |V_{tb}|^2 \approx 1.0
            \quad (\text{kinematically suppressed by }m_t \approx 173\,\mathrm{GeV}).
\end{align}
Off-diagonal CKM contributions ($d\to c$, $s\to u$, etc.) are suppressed by
the Cabibbo angle and are included by POWHEG but are numerically subdominant.
Antiquark sea contributions, including $\bar{u}\to \bar{d}$ and
$\bar{c}\to \bar{s}$ transitions, also enter through the PDF at small $x$.
At NLO, the gluon PDF contributes through the splitting subprocess
$\nu_\ell + g\to \ell^- + q_j + \bar{q}_i$. The resulting outgoing
parton flavors that can appear in the LHE hard-process record are
$u$, $c$, $d$, $s$, and $g$, with the exact flavor mix determined by
PDF support and POWHEG matrix-element channel selection at the given $(x,Q^2)$.

\section{LHE hard-process event structure}

The Les Houches Event (LHE) file written by \texttt{pwhg\_main} contains one
\texttt{<event>} block per generated hard-process realization. Each block
records the incoming and outgoing particles with their PDG identifiers,
four-momenta $(p_x,p_y,p_z,E,m)$ in GeV, and a status code: $-1$ for
incoming, $+1$ for outgoing. The event header also carries the POWHEG event
weight and the hardest-emission scale $\mu_{\rm POWHEG}$ that will serve as
the starting scale for the subsequent parton shower.

At Born level the event topology is
\begin{equation}
\begin{array}{ll}
\text{status}=-1: & \nu_\ell + q_i,\\
\text{status}=+1: & \ell^- + q_j.
\end{array}
\end{equation}
With the NLO real emission included, the topologies are
\begin{equation}
\begin{array}{ll}
\text{status}=-1: & \nu_\ell + q_i,\\
\text{status}=+1: & \ell^- + q_j + g
\end{array}
\qquad\text{or}\qquad
\begin{array}{ll}
\text{status}=-1: & \nu_\ell + g,\\
\text{status}=+1: & \ell^- + q_j + \bar{q}_i.
\end{array}
\end{equation}
The interpretation of the event weight depends on the LHE \code{IDWTUP}
strategy and is discussed in the provenance and statistical-weights chapter;
the raw sum of \code{XWGTUP} values is not, in general, the total cross section.

The LHE particles are parton-level hard-process records. In H3-W10 they are
the immutable input to PYTHIA~8 showering, hadronization, and decays. GEANT4
transport remains disabled; therefore the H3-W10 output is a generator-level
hadronic final state, not a detector-level observable.

\section{Dry-run role}

H3-W9a prepares local POWHEG DIS jobs but deliberately does not execute
\code{pwhg_main}. It consumes primary Observer Image Branch records and writes its own
products under \code{POWHEG/}. It does not modify \code{ObserverBridge/},
\code{ObserverImageBranches/},
\code{DIS/}, \code{ForwardGeodesics/}, or \code{UHEsource/}.

The H3-W9a run mode is
\begin{equation}
\code{powheg_run_mode}=\code{dry_run}.
\end{equation}
The reports record
\begin{align}
\code{powheg_invoked} &= \code{false},\\
\code{pwhg_main_executed} &= \code{false},\\
\code{powheg_lhe_generated} &= \code{false}.
\end{align}

H3-W9b adds a controlled real-smoke mode:
\begin{equation}
\code{powheg_run_mode}=\code{real_smoke}.
\end{equation}
This mode executes only the first primary branch record, writes one input card, executes
the self-contained local executable
\code{external/powheg/build/DIS/pwhg_main}, and validates a minimal LHE file.
It is a pipeline smoke execution, not a precision production campaign.
Perturbative order is independent of run mode: both LO and NLO smoke cards are
supported and are labeled in summaries and provenance. In H3-W9b itself,
PYTHIA, GEANT4, photon transport, and spectra remain disabled; H3-W10 consumes
the resulting validated LHE in a separate output directory.

\section{POWHEG input branch contract}

The driver reads
\code{ObserverImageBranches/observer_image_primary_branches.jsonl}. The
selection of downstream candidates is performed by the Observer Bridge and the
selection of the primary observed image is performed by H3-W8b. POWHEG does not
rank candidates and does not choose among multiple image branches. It prepares
one POWHEG request for each primary branch record, except in real-smoke mode
where the Python wrapper intentionally executes only the first prepared request
as a safety smoke test.

\section{Input-card mapping}

Each selected candidate receives a request identifier of the form
\code{H3PWHG-000001}. Each request has its own
\code{powheg_input_cards/<request_id>/powheg.input}. The implemented mapping is
\begin{align}
E_{\nu,\rm local} &\rightarrow \code{ebeam1},\\
m_N = 0.938272~\GeV &\rightarrow \code{ebeam2},\\
\code{events_per_candidate} &\rightarrow \code{numevts},\\
\code{fixed_target} &= 1,\\
\code{masslesslhe} &= 1,\\
\code{channel_type} &= 3,\\
\code{perturbative_order=LO} &\rightarrow \code{LOevents=1},\\
\code{perturbative_order=NLO} &\rightarrow \text{no \code{LOevents} flag}.
\end{align}
The dry-run card also records fixed PDF choices
\code{lhans1 = 303400} and \code{lhans2 = 303400}. The implemented scale cap is
\begin{equation}
Q_{\max} =
\min\left(\sqrt{2(0.938272)\,E_{\nu,\rm local}},\,10^5\right).
\end{equation}
For the implemented seed mode, the seed is
\begin{equation}
\code{powheg_seed}
=
\code{random_seed}+\code{candidate_rank}.
\end{equation}

\section{Real-smoke and real-free LHE validation}

In H3-W9b real-smoke mode, the prepared POWHEG card is copied into an isolated
work directory and \code{pwhg_main} is executed with that directory as its
working directory. The produced \code{pwgevents.lhe} must satisfy
\begin{align}
\code{<LesHouchesEvents} &\in \code{pwgevents.lhe},\\
\code{</LesHouchesEvents>} &\in \code{pwgevents.lhe},\\
N_{\rm LHE\,events} &> 0.
\end{align}
The validated LHE and run log are copied to
\begin{align}
\code{POWHEG/powheg_lhe/H3PWHG-000001/pwgevents.lhe},\\
\code{POWHEG/powheg_run_logs/H3PWHG-000001/powheg.log}.
\end{align}
The validation report records \code{powheg_invoked=true},
\code{pwhg_main_executed=true}, \code{powheg_lhe_generated=true}, and
\code{n_lhe_events}. This smoke mode does not modify
\code{ObserverImageBranches/}, \code{ObserverBridge/}, \code{DIS/}, \code{ForwardGeodesics/}, or
\code{UHEsource/}.

H3-W9b real-free mode uses the same local POWHEG execution and LHE validation
rules, but it removes the real-smoke safety clamp. The number of interaction
candidates is set by the number of records in
\code{ObserverImageBranches/observer_image_primary_branches.jsonl}, and the
number of independent POWHEG hard-scattering realizations per selected
interaction is set by \code{events_per_candidate}. Each primary branch record is executed in its own
isolated POWHEG work directory. The produced LHE files are copied under
\code{POWHEG/powheg_lhe/<powheg_request_id>/pwgevents.lhe}, and the parser
aggregates all produced LHE files into the POWHEG particle and event summaries.
The provenance records \code{powheg_real_free_invoked=true},
\code{real_free_mode=true}, \code{n_powheg_jobs_requested},
\code{n_powheg_jobs_run}, \code{events_per_candidate_requested}, and
\code{n_lhe_events_total}. This mode itself does not invoke PYTHIA, GEANT4,
photon transport, or spectra; PYTHIA is invoked only by H3-W10.

\section{Event requests}

The event request file
\code{POWHEG/powheg_event_requests.jsonl} records for each prepared job:
\begin{align}
\mathcal{R}_{\rm POWHEG}=\{&
\code{powheg_request_id},\,
\code{candidate_rank},\,
\code{interaction_id},\,
\code{event_id},\,
\code{source_sample_id}, \nonumber\\
&E_{\nu,\rm local},\,
w_{\rm physics},\,
w_{\rm observer},\,
S_{\rm obs},\,
\code{powheg_input_path},\,
\code{powheg_seed},\,
\code{powheg_status}\}.
\end{align}
The dry-run status is \code{dry_run_ready}; \code{powheg_invoked} is false for
every dry-run request. In real-smoke mode the initial status is
\code{real_smoke_ready}; after successful validation the request records
\code{real_smoke_lhe_generated}, \code{powheg_invoked=true},
\code{pwhg_main_executed=true}, the LHE path, the log path, and
\code{n_lhe_events}. In real-free mode the analogous statuses are
\code{real_free_ready} and \code{real_free_lhe_generated}; multiple validated
LHE products may be recorded in a single POWHEG run.

\section{Inputs and outputs}

\begin{longtable}{p{0.30\textwidth}p{0.60\textwidth}}
\toprule
Role & Implemented item \\
\midrule
Inputs & \code{ObserverImageBranches/observer_image_primary_branches.jsonl} and POWHEG
configuration values.\\
Outputs & \code{POWHEG/powheg_event_requests.jsonl}, input cards, summary
JSON/CSV, report JSON, validation report, diagnostic plots, and in real-smoke
or real-free mode validated \code{pwgevents.lhe} products.\\
Still disabled in H3-W9 & No PYTHIA inside the POWHEG stage; no GEANT4;
no photon transport; no spectra.\\
Backend & C++17 executable \code{bin/hadros3_powheg_driver}. Python handles
orchestration and plots.\\
\bottomrule
\end{longtable}

\chapter{H3-W10 Event Generation with PYTHIA 8 and HepMC3}

\section{Boundary and execution modes}

H3-W10 consumes only validated LHE files produced by the H3-W9b
\code{real_smoke} or \code{real_free} modes. A dry-run card is not an event and
is rejected. The implemented transformation is
\begin{equation}
\mathrm{POWHEG\ LHE}\longrightarrow
\mathrm{PYTHIA\ 8\ shower/hadronization/decays}\longrightarrow
\mathrm{HepMC3}+\mathrm{JSONL}.
\end{equation}
The \code{dry_run} mode validates discovery and writes a manifest without
calling PYTHIA; \code{parton_check} parses the LHE and preserves its final
partons without a shower; \code{real_smoke} clamps execution to one request
and at most two events; and \code{real_free} uses the configured positive
limits. GEANT4, photon transport, and detector response remain disabled.

Every accepted LHE event maps one-to-one to an event identifier
\begin{equation}
\code{event_generation_event_id}
=\code{powheg_request_id}:\code{lhe_event_index}.
\end{equation}
No candidate resampling or implicit oversampling occurs in this stage.

\section{Fixed-target local frame}

The generator frame is the local matter tetrad used to prepare the fixed-target
POWHEG card. At ultra-high energy the POWHEG laboratory reboost can leave the
incoming lepton energy slightly different from the LHE \code{EBMUP}. Before
external showering, H3-W10 applies one common longitudinal Lorentz boost to
every particle of an event,
\begin{align}
E'   &= \cosh y\,E-\sinh y\,p_z,\\
p_z' &= \cosh y\,p_z-\sinh y\,E,\\
y    &= \ln(E_{\ell,\mathrm{event}}/E_{\ell,\mathrm{beam}}).
\end{align}
This maps the incoming massless lepton to the declared beam energy while
preserving invariant masses and event topology. After PYTHIA, the complete
event is transformed back to the target local energy. Conservation is tested
against the physical fixed-target initial state
\begin{equation}
P^\mu_{\rm initial}=(E_\nu+m_p,\,0,\,0,\,E_\nu),
\qquad m_p=0.938272~\GeV,
\end{equation}
with momenta in GeV and HepMC vertices in millimetres.

\section{Shower, matching, hadronization, and decays}

PYTHIA reads the LHE through \code{Beams:frameType=4}; lepton PDFs are disabled
with \code{PDF:lepton=off}, multiple parton interactions are disabled, and the
POWHEG hook vetoes shower emissions above the LHE \code{SCALUP}. The accepted
H3-W10 policy enables final-state radiation, hadronization, and decays, but
disables generic initial-state radiation.

This ISR restriction is numerical rather than cosmetic. In the integrated UHE
fixed-target smoke, generic PYTHIA ISR produced a relative four-momentum
residual of $2.61\times10^{-5}$, which fails the H3-W10 tolerance. With ISR
disabled and POWHEG-vetoed FSR, hadronization, and decays enabled, the measured
residual was $9.47\times10^{-10}$. ISR therefore remains experimental until a
DIS fixed-target hook passes the same conservation test. The effective PYTHIA
settings, versions, compiler, backend hash, seed, matching policy, and LHE hash
are recorded in the manifest and per-job \code{pythia.cmnd} dump.

For each generated event H3-W10 requires
\begin{align}
\delta_P &=
\frac{\lVert P_{\rm final}-P_{\rm initial}\rVert_2}
{\max(E_{\rm initial},1~\GeV)} \leq 5\times10^{-8},\\
\delta_{m^2} &\leq 10^{-8},\\
Q_{\rm shower,max} &\leq \code{SCALUP}
\quad\text{within the backend numerical allowance}.
\end{align}
Charge conservation is audited, the final hadronized record must contain
hadrons, and no final-state quarks or gluons may remain when hadronization is
enabled.

\section{Statistical weights and products}

The signed LHE \code{XWGTUP} value is preserved event by event. Its
interpretation follows \code{IDWTUP}; in particular, H3-W10 does not treat a
negative NLO weight as a selection probability and does not assume that the
raw weight sum is a cross section. The fields \code{physics_weight},
\code{observer_weight}, and \code{final_observation_score} remain separate
metadata. No multiplication of generator and observer weights is performed.

Official products are written only under \code{EventGeneration/}: immutable
input hashes and configuration in \code{event_generation_manifest.json},
per-job records, aggregate HepMC3, event and final-particle JSONL, validation
and particle-content reports, summary JSON/CSV, four diagnostics, and a local
event viewer. Rerunning POWHEG invalidates these downstream products, while an
H3-W10 run verifies that the upstream LHE hashes did not change.

\chapter{Provenance and Statistical Weights}

\section{Stage provenance}

Each stage records the backend, products, validation flags, and disabled future
stages relevant to its current implementation. A typical provenance record
contains:
\begin{equation}
\{\mathrm{parameters},\mathrm{products},\mathrm{validation},
\mathrm{stage\ summaries},\mathrm{disabled\ stages}\}.
\end{equation}
The implemented provenance reflects the selected execution mode. In a real
H3-W10 run it records \code{pythia_invoked=true} and
\code{expensive_event_generation_invoked=true}; dry-run and parton-check modes
leave both false. Later transport stages remain disabled:
\begin{align}
\code{geant4_invoked} &= \code{false},\\
\code{photon_transport_invoked} &= \code{false},\\
\code{spectra_invoked} &= \code{false}.
\end{align}
For H3-W9a, POWHEG is active only as dry-run preparation:
\begin{equation}
\code{status}=\code{powheg_jobs_prepared_no_lhe}.
\end{equation}
For H3-W9b real-smoke execution:
\begin{equation}
\code{status}=\code{powheg_real_smoke_lhe_generated}.
\end{equation}
The provenance also records the scientific theory reference for the run:
\begin{align}
\code{theory_document} &= \code{docs/Theory/HADROS3_Physics_Theory.pdf},\\
\code{theory_version} &= \code{\TheoryVersion},\\
\code{theory_commit} &= \code{runtime git commit},\\
\code{theory_generation_date} &= \code{runtime provenance timestamp}.
\end{align}
These fields connect each generated result to the theory document that describes
the implemented physical and statistical assumptions.

\section{Weight chain}

The implemented source and direction samplers are configured so that the
default source and direction weights are unity:
\begin{align}
w_{\rm source} &= 1,\\
w_{\rm direction} &= 1.
\end{align}
The DIS sampler records interaction and pre-event weights. The Observer Bridge
then uses
\begin{equation}
w_{\rm physics}=\code{final_pre_event_weight}
\end{equation}
and computes
\begin{equation}
S_{\rm obs}=w_{\rm physics}w_{\rm observer}.
\end{equation}
The POWHEG dry-run and real-smoke modes carry those weights into each event
request. They do not change the physics score. H3-W9a stops after card
preparation; H3-W9b executes one local POWHEG smoke job and validates the
resulting minimal LHE.

\section{LHE weight normalization}

The parser reads the complete LHE initialization header, including
\code{IDWTUP}, and each process row
\begin{equation}
(\code{XSECUP},\code{XERRUP},\code{XMAXUP},\code{LPRUP}).
\end{equation}
The generator-declared total cross section is
\begin{equation}
\sigma_{\rm init}=\sum_{p=1}^{\code{NPRUP}}\code{XSECUP}_p.
\end{equation}
The current POWHEG output uses \code{IDWTUP=-4}. For this weighted-event
strategy the event-sample estimator recorded by HADROS3 is the signed mean
\begin{equation}
\widehat{\sigma}_{\rm event}
=\frac{1}{N}\sum_{i=1}^{N}\code{XWGTUP}_i,
\end{equation}
with the usual standard error of the sample mean. The raw sum is stored only as
a diagnostic and is explicitly marked as not being a cross-section estimator.
This distinction makes the estimate invariant if the same event sample is
duplicated. For other \code{IDWTUP} values, HADROS3 reports
$\sigma_{\rm init}$ and requires the corresponding generator-specific
normalization rule rather than silently applying the \code{IDWTUP=-4} formula.

\section{Random seeds}

The implemented random stages record seeds:
\begin{itemize}
\item UHE source sampling records its source seed.
\item DIS interaction sampling records its acceptance and point-sampling seed.
\item POWHEG dry-run and real-smoke modes record \code{random_seed} and the
per-candidate \code{powheg_seed}.
\end{itemize}
These seeds are part of the statistical contract of the pipeline.

\chapter{H3-W11 GEANT4 Local Material Transport}

\section{Scientific boundary}

H3-W11 consumes the H3-W10 HepMC3 final state and transports supported
particles through a finite material patch in the orthonormal local matter
tetrad.  The patch origin is the DIS interaction vertex.  Its Cartesian axes
are the spatial legs of the local tetrad; energy and momentum use GeV, while
HepMC3 and GEANT4 positions use mm.  H3-W11 does not integrate Kerr geodesics.
The transformation from the local tetrad to the global Kerr frame remains
metadata for H3-W12.

Each accepted HepMC3 event maps to one \code{G4Event}.  Final particles map to
\code{G4PrimaryParticle} objects without resampling.  Generator, physics, and
observer weights remain separate from the initially unit GEANT4 statistical
track weight.  The official output records local boundary crossings, track
and parent identifiers, creator process, position, four-momentum, local time,
and boundary normal for later H3-W12 propagation.

\section{Geometry and material}

The implemented geometry is \code{local_homogeneous_patch_v1}: a box centered
on the interaction.  Vacuum surrounds the patch and all six boundary faces
are audited.  In official runs the density is recovered by
\code{interaction_id} from the originating H3-W7 accepted interaction.  A
configured density is permitted only in explicitly labelled validation
fixtures.  The baseline composition is a configurable H/He mixture, with
mass fractions and density written to the manifest.  This is a controlled
local-medium prescription, not a laboratory detector and not a full GRMHD
toroidal geometry.

The baseline includes no magnetic or electric field and no optical-photon
physics.  Neutrinos use transport-only pass-through behavior unless an
explicitly validated interaction model is later registered.

\section{Physics list and validity domain}

The implemented reference list is \code{FTFP_BERT} from GEANT4 11.4.2.  The
H3-W11 v1 contract imposes the conservative upper bound
\begin{equation}
 E_{\rm primary} \leq 10^5\ {\rm GeV}.
\end{equation}
This bound follows the documented upper range of the FTF component for the
covered hadrons and must not be interpreted as validation above that range.
Before \code{BeamOn}, every final primary is classified by species and energy.
An out-of-domain particle produces \code{unsupported_domain}, writes an audit,
and prevents transport.  Energy clipping and silent extrapolation are
forbidden.

The H3-W10 event audited on 2026-07-13 contained 57 final primaries, of which
22 exceeded the bound.  Its largest primary energy was
\(3.3973576244\times10^8\) GeV.  Consequently the correct H3-W11 result for
that event is a domain refusal with zero transported events, rather than an
unvalidated GEANT4 shower.

\section{Scoring and energy ledger}

For each event the backend records initial total primary energy, local energy
deposit, total energy crossing the patch boundary, steps, and a raw balance.
In vacuum the unexplained normalized residual must be below \(10^{-6}\).
Inside matter, rest-mass creation, annihilation, nuclear binding, and the
implicit material reservoir prevent the naive identity
\(E_{\rm in}=E_{\rm dep}+E_{\rm escape}\) from being used as a precision
conservation test.  H3-W11 therefore records the raw balance separately as an
inferred material rest-mass and binding exchange; it is never hidden by
renaming the raw quantity.

Numerical validation fixtures test six-axis vacuum crossings, pion decay,
photon attenuation in lead, muon stopping-power scaling in water, proton
interaction-length scaling in lead, neutrino pass-through, deterministic
seeds, signed-weight preservation, atomic publication, and the UHE domain
guard.  Fixture results validate the implementation inside its declared
domain; they do not validate production UHE transport.

\section{Execution and products}

Opening or reloading the GEANT4 web tab never starts transport.  Execution
requires an explicit POST action or a command-line mode.  The official
\code{GEANT4/} directory contains the summary, validation report, environment
and physics manifest, import/domain audit, transported event summaries,
escaping-particle JSONL, unsupported-particle JSONL, plots, logs, and event
viewer.  Publication is staged and atomic.  A new H3-W10 product invalidates
H3-W11 and all later transport products.

The precise implementation status is therefore
\begin{quote}
H3-W11 implemented and numerically tested for the declared supported-energy
domain; real HADROS3 UHE production remains blocked by the physics-domain
guard.
\end{quote}

\chapter{Pipeline Summary}

\section{Implemented stage contract}

The current implemented pipeline can be summarized as
\begin{equation}
\mathrm{UHESource}
\rightarrow
\mathrm{ForwardGeodesics}
\rightarrow
\mathrm{DIS}
\rightarrow
\mathrm{ObserverBridge}
\rightarrow
\mathrm{POWHEG\ RealSmoke}
\rightarrow
\mathrm{EventGeneration}
\rightarrow
\mathrm{GEANT4}_{\rm supported\ energy}.
\end{equation}
The camera preview is an observer-view diagnostic that can be used by the
dashboard and by Observer Bridge overlays, but it is not a physical
secondary-particle propagation stage.

\begin{longtable}{p{0.20\textwidth}p{0.24\textwidth}p{0.24\textwidth}p{0.22\textwidth}}
\caption{Implemented pipeline products and semantic status.}\\
\toprule
Stage & Official inputs & Official outputs & Status \\
\midrule
\endfirsthead
\toprule
Stage & Official inputs & Official outputs & Status \\
\midrule
\endhead
UHE Source & Configuration & \code{UHEsource/} JSONL/JSON & Sampling proxy with
monoenergetic source.\\
Forward Geodesics & \code{UHEsource/} & \code{ForwardGeodesics/} paths and
segments & Kerr null geodesic propagation.\\
DIS Sampler & \code{ForwardGeodesics/}, \code{UHEsource/} & \code{DIS/}
optical-depth and accepted interactions & Optical-depth DIS sampler.\\
Observer Bridge & \code{DIS/dis_accepted_interactions.jsonl} &
\code{ObserverBridge/} candidates and ranked events & Non-destructive
geometric scoring proxy.\\
Gravitational Image Analysis & \code{ObserverBridge/} selected candidates and Kerr
pixel-map products & \code{ObserverImageBranches/} branches and primary
branches & Observed image branch proxy selection.\\
POWHEG Real Smoke & \code{ObserverImageBranches/observer_image_primary_branches.jsonl} &
\code{POWHEG/} requests, input cards, validation report, and one smoke LHE &
H3-W9a dry run plus H3-W9b local one-job LHE smoke.\\
Event Generation & Validated real H3-W9b LHE files & \code{EventGeneration/}
HepMC3, JSONL, reports, plots, and manifest & H3-W10 PYTHIA~8 FSR,
hadronization, decays, and HepMC3 conversion; ISR disabled by numerical policy.\\
GEANT4 & Validated H3-W10 HepMC3 and H3-W7 local density metadata &
\code{GEANT4/} audits, manifests, boundary particles, ledgers, plots, and logs &
H3-W11 local material transport inside the declared supported-energy domain;
UHE inputs are blocked before transport.\\
\bottomrule
\end{longtable}

\section{Explicit non-goals of the implemented stages}

The current pipeline does not yet implement:
\begin{itemize}
\item an accepted fixed-target DIS initial-state shower; generic PYTHIA ISR is
experimental and disabled by default after failing the conservation tolerance;
\item GEANT4 transport of the current UHE event beyond the declared
supported-energy envelope;
\item photon transport from secondaries to the observer;
\item full GRMHD torus evolution;
\item full physical visibility, escape, redshift, or line-of-sight transport.
\end{itemize}
Those future stages must preserve the contract that each stage consumes official
previous outputs, writes only to its own directory, records provenance, and
marks expensive or proxy physics explicitly.

\appendix
\chapter{Appendix: Symbols and Product Names}

\begin{longtable}{p{0.30\textwidth}p{0.60\textwidth}}
\toprule
Product or field & Meaning \\
\midrule
\endfirsthead
\toprule
Product or field & Meaning \\
\midrule
\endhead
\code{uhe_neutrino_source_samples.jsonl} & UHE source positions, energies,
directions, and sampling weights.\\
\code{uhe_neutrino_forward_path_segments.jsonl} & Forward geodesic path
segments with midpoint coordinates, momenta, and segment lengths.\\
\code{dis_path_optical_depths.jsonl} & Path-level DIS optical depths and
interaction probabilities.\\
\code{dis_accepted_interactions.jsonl} & Accepted DIS interactions and their
density-checked interaction points.\\
\code{observer_bridge_candidates.jsonl} & All accepted DIS interactions scored
as observation candidates.\\
\code{observer_bridge_ranked_events.jsonl} & Same candidate family ranked by
final observation score.\\
\code{observer_image_primary_branches.jsonl} & Primary observed image branch
per selected Observer Bridge candidate.\\
\code{powheg_event_requests.jsonl} & Dry-run POWHEG event requests and card
paths.\\
\code{final_pre_event_weight} & Physics weight consumed by Observer Bridge as
implemented.\\
\code{final_observation_score} & Product of implemented physics and observer
weights.\\
\code{powheg_invoked} & False in H3-W9a dry run; true in H3-W9b real smoke.\\
\code{pwhg_main_executed} & False in H3-W9a dry run; true in H3-W9b real smoke.\\
\code{powheg_lhe_generated} & False in H3-W9a dry run; true in H3-W9b real smoke.\\
\code{event_generation_events.hepmc3} & Aggregate H3-W10 generator-level event
record in HepMC3 ASCII format.\\
\code{xwgtup} & Signed LHE generator weight, kept separate from observer scores.\\
\code{pythia_invoked} & True only in H3-W10 \code{real_smoke} or
\code{real_free}; false in dry-run and parton-check modes.\\
\code{geant4_escaped_particles.jsonl} & H3-W11 local boundary crossings ready
for H3-W12, including lineage, process, local phase space, time, and weight.\\
\code{geant4_invoked} & True only when \code{BeamOn} transported an input fully
inside the declared domain; false for import checks and UHE refusals.\\
\bottomrule
\end{longtable}

\end{document}
